\documentclass[10pt,a4paper]{amsart}
\usepackage[margin=19mm]{geometry}
\usepackage{amsmath,amssymb,mathtools}
\usepackage[scr=rsfs,cal=euler]{mathalfa}
\usepackage{xfrac}
\usepackage[unicode,hidelinks,bookmarks=false]{hyperref}
\newcommand{\ie}{i.e.\ }
\newcommand{\cf}{cf.\ }
\newcommand{\resp}{respectively\ }
\newcommand{\Subsection}{\subsection}
\providecommand{\nobreakdash}{\nobreak\hbox{-}}
\setlength{\emergencystretch}{2em}
\multlinegap0pt
\usepackage{dsfont}
\usepackage{amssymb}
\usepackage{enumerate}
\newtheorem{lemma}{Lemma}
\title{Unbounded mean convex domains in Euclidean space}
\author{Jian Ge}
\date{Source: arXiv:2605.16802v1 (CC BY 4.0)}
\begin{document}
\maketitle

\noindent\emph{Source and licence.} Unbounded mean convex domains in Euclidean space. Version arXiv:2605.16802v1 (CC BY 4.0), distributed by arXiv under the Creative Commons Attribution 4.0 International licence, http://creativecommons.org/licenses/by/4.0/, as recorded in the arXiv licence metadata for this exact version and shown on the abstract page https://arxiv.org/abs/2605.16802v1. Full licence notice: \texttt{licenses/arxiv-2605.16802-CC-BY-4.0.txt}.\par\smallskip

\noindent\textit{Editorial scope.} Counted unit: the article's lemma lem:Hessofr (source0.tex lines 251--261). The article's complete original proof of it is delivered with it (source0.tex lines 262--277, one \texttt{proof} environment of 16 lines). No mathematics is written, repaired or re-derived: the statement and the proof are the article's own lines, in the article's own notation, and no retained line is substituted or re-flowed.  Retained uncounted as the source-local context the proof needs: lines 241--243.  Nothing was omitted from any retained span, so the package carries no omission record.  No retained line contains a citation, so the wrapper carries no bibliography and no bibliography entry.  No retained line contains a figure environment, an image-inclusion command or drawing code, so no image is delivered and the package declares no asset dependency.  Editorial formatting only, in addition: an AMS \texttt{amsart} preamble that restates the article's own declarations and macros used by the retained lines, namely the environments \texttt{lemma} on separate counters, together with the standard packages those lines need.  The retained environments print in this document as Lemma 1 in order of appearance, whereas the article prints the same items with the numbering of its own full document.  The complete native source of the article as distributed in the arXiv e-print for this exact version is bundled unchanged as \texttt{sources/200747/source0.tex}.
\begin{equation}\label{eq:KEY}
	|v_{i}-w_{i}| \le C \varepsilon\sqrt{R}\to 0.
\end{equation}

\begin{lemma}\label{lem:Hessofr}
	Let notations be as above, we have:
	\[
		\operatorname{Hess}_{\Sigma\times\Gamma}r(V_i,V_i) = Q_{i} - a h_\Sigma(v_i, v_i)-b h_{B}(w_i, w_i).
	\]
	where
	\begin{equation}\label{eq:KEY01}
		0\le Q_{i}:= \frac{|v_{i}-w_{i}|^{2} - \left\langle v_{i}-w_{i}, \uparrow_{x}^{y} \right\rangle ^{2}}{R}\le C \varepsilon^{2},
	\end{equation}
	for some constant $C>0$ independend of $\varepsilon$.
\end{lemma}

\begin{proof}
	Let $\alpha_{i}(s)\subset \Sigma$ and $\beta_{i}(s)\subset B$ be geodesics in $\Sigma$ and $B$ respectively such that
	\begin{equation*}
	\alpha_{i}(0)=x, \alpha_{i}'(0)=v_{i}; \quad \beta_{i}(0)=y, \beta_{i}'(0)=w_{i}.
	\end{equation*}
	Then
	\begin{equation}\label{eq:temp02}
	\alpha_{i}''(0)=-h_{\Sigma}(v_{i}, v_{i}) \nu_{\Sigma}(x);\quad \beta_{i}''(0)=-h_{B}(w_{i}, w_{i})\nu_{B}(y).
	\end{equation}
	Here $h_{\Sigma}$ is the second fundamental form of $\Sigma$ and $\nu$ is the outward normal of $\partial X$. Standard Euclidean calculation shows.
	\begin{equation}\label{eq:temp03}
		r''(0) = \frac{|v_{i}-w_{i}|^{2} - \left\langle v_{i}-w_{i}, \uparrow_{x}^{y} \right\rangle ^{2}}{r} 
		- \left\langle \uparrow_{x}^{y}, \alpha''(0)-\beta''(0) \right\rangle.
	\end{equation}
	Combine \eqref{eq:temp02} and \eqref{eq:temp03}, we have the desired Hessian form. The estimate of $Q$ follows from \eqref{eq:KEY}.
\end{proof}

\end{document}
